% ============================================================
% 3. WHY FIXED-DIRECTION GATING FAILS
% ============================================================
\section{Why Fixed-Direction Gating Fails}
\label{sec:signal-landscape}
\begin{table}[!t]
\caption{Signal--utility correlations across 6 environments and
3 backbones. $\rho$: Spearman correlation of token entropy with
optimizer utility $U$. $^*$: significant at $p{<}0.05$.
Direction varies across \emph{both} environments and model
backbones.}
\vspace{-0.1in}
\label{tab:signal-discovery}
\centering\small
\setlength{\tabcolsep}{4pt}
\begin{tabular}{lcccc}
\toprule
Environment & Qwen3 & Phi-3.5 & Llama-3.1 & Cons. \\
\midrule
HotpotQA   & $-$0.327$^*$ & $+$0.184$^*$ & $-$0.346$^*$ & \xmark \\
WebShop    & $+$0.133$^*$ & $+$0.335$^*$ & $+$0.287$^*$ & \cmark \\
FEVER      & $-$0.119$^*$ & $-$0.156$^*$ & $+$0.428$^*$ & \xmark \\
APPS       & $+$0.317$^*$ & $-$0.024     & $-$0.249$^*$ & \xmark \\
TWExpress  & $-$0.290$^*$ & $+$0.000     & $+$0.000     & --- \\
Plancraft  & $-$0.016     & $+$0.000     & $-$0.176$^*$ & --- \\
\bottomrule
\end{tabular}
\end{table}

We first formalize the adaptive gating problem, then measure
signal--utility correlations across diverse environments to test
the fixed-direction assumption.

We study LLM-based agents in multi-step interactive environments.
At each step $t$, the agent observes state $s_t$, selects action
$a_t$, and transitions to $s_{t+1}$. An environment-specific
\emph{test-time optimizer} $T$ (e.g., rollout evaluation, $K$-variant
sampling) can be invoked at any step to potentially improve action
quality. Let $\sigma(s)$ denote the uncertainty signal at state $s$
(e.g., token entropy of the next-token distribution). We define
the \emph{optimizer utility} as:
\begin{equation}
  U(T, s) = \mathbb{E}[R(\tau) \mid a{=}T(s)]
           - \mathbb{E}[R(\tau) \mid a{=}\pi(s)]
  \label{eq:utility}
\end{equation}
where $R(\tau)$ is the episode return. $U > 0$ indicates the
optimizer improves the outcome; $U \leq 0$ indicates it is wasteful
or harmful. The \emph{adaptive gating problem} is to learn a gate
$g: \mathcal{S} \to \{0, 1\}$ that triggers $T$ only when $U > 0$.

\subsection{Signal--Utility Direction Reverses Across Environments
            and Backbones}
\label{sec:empirical-landscape}

\textbf{The signal--utility direction is not fixed.}
Table~\ref{tab:signal-discovery} reports Spearman correlations
between token entropy and optimizer utility across 6 environments
<<<<<<< HEAD
and 3 model backbones~\cite{DBLP:journals/corr/abs-2505-09388,DBLP:journals/corr/abs-2407-21783,DBLP:journals/corr/abs-2404-14219}. Two patterns stand out, and the more striking one is the harder to explain. First, \emph{the same signal in the same environment can reverse across
backbones}: in 3 out of 4 environments with multiple significant
correlations, the backbones disagree on the sign of $\rho$. The
most extreme case is FEVER, which flips from $\rho{=}{-}0.156$ on
Phi-3.5 to $\rho{=}{+}0.428$ on Llama-3.1---a sign change driven
=======
and 3 model backbones~\cite{DBLP:journals/corr/abs-2505-09388,DBLP:journals/corr/abs-2407-21783,DBLP:journals/corr/abs-2404-14219}. First, \emph{the same signal in the same environment can reverse across
backbones}: among environments with at least two significant nonzero correlations, sign disagreement across backbones is common. The most extreme case is FEVER, which flips from $\rho{=}{-}0.156$ on Phi-3.5 to $\rho{=}{+}0.428$ on Llama-3.1, a sign change driven
>>>>>>> 1bf5b0e2c05afb339eb3e76f610bbf4ff034c5a3
entirely by the choice of LLM, with the task held fixed. The
signal's semantics are therefore not a property of the
environment alone, but of how the agent's prior knowledge
interacts with it. Second, \emph{even within a single backbone,
the direction also varies across environments}: on Qwen3, entropy
is negatively correlated with utility in information-seeking
environments (FEVER, HotpotQA, TWExpress) but positively in code
generation (APPS) and web navigation (WebShop), with near-zero
correlation in weak-signal environments (Plancraft). Crucially,
this direction mismatch is not benign: as
Section~\ref{sec:main-results} will show, gates that assume the
wrong direction systematically fall \emph{below} the no-trigger
baseline, making them strictly worse than doing nothing in
mismatched (environment, backbone) pairs.
<<<<<<< HEAD

\paragraph{Implication: $\sigma$ is not a sufficient statistic
for VOC.}
The reversal in Table~\ref{tab:signal-discovery} is not just
``fixed-direction gating fails''---it is the empirical signature
of a stronger structural fact. Across (environment, backbone)
settings, two states with identical $\sigma$ values can have
opposite-sign value of computation, depending on whether the
agent is in an information-poor or decision-difficult state
(\S\ref{sec:toy-model}). Hence $\sigma$ is not a sufficient
statistic for $\mathrm{VOC}$, and \emph{no} policy of the form
$g(\sigma(s))$---monotone or otherwise, calibrated or
not---can deliver non-negative value of computation across
heterogeneous settings (Proposition~\ref{prop:necessity}, proof
in Appendix~\ref{app:proof-necessity}). Fixed-direction gating
is the strict-monotone special case; the underlying obstruction
is a property of the signal itself, not of how the gate
post-processes it. Recovering sufficiency therefore requires
augmenting $\sigma$ with auxiliary features that distinguish
the two state types---a requirement that directly motivates
the multi-feature design of \S\ref{sec:method}.

\paragraph{A diagnostic for practitioners.}
The reversal test of Table~\ref{tab:signal-discovery} is itself
the leave-behind tool of this paper: collecting a few hundred
signal-agnostic $(s, U)$ pairs and reporting the Spearman
correlation $\rho(\sigma, U)$ together with its bootstrap
confidence interval is sufficient to decide whether a fixed
direction is safe for a given (environment, backbone) deployment.
A confidence interval that crosses zero, or two settings that
disagree on sign, indicates that any uncertainty-thresholding
gate will eventually fall below never-gating; a stable
single-sign correlation indicates that fixed-direction gating is
safe and direction learning is unnecessary. We apply this
diagnostic throughout, and recommend it as a standalone check
independent of the specific learning procedure adopted.
=======
>>>>>>> 1bf5b0e2c05afb339eb3e76f610bbf4ff034c5a3

\textbf{Signal identity also varies across environments.}
Beyond direction, the most informative signal itself differs
across environments: \texttt{step\_count} dominates in most
settings, but WebShop relies on \texttt{num\_available\_actions}
instead (Appendix~\ref{app:feature-selection}). No single signal
is reliably informative; multi-signal gates substantially
outperform single-feature gates
(Appendix~\ref{app:auc}). These observations suggest that
uncertainty is entangled with latent state structure that varies
across (environment, backbone) pairs. We formalize this intuition
next.

\subsection{Why Does Direction Reverse? A Two-Source Model}
\label{sec:toy-model}

The findings above establish that fixed-direction gating fails; we
<<<<<<< HEAD
now ask \emph{why}. We trace direction reversal to two coexisting
state types with opposite signal--utility relationships.

\paragraph{Intuition.}
Consider two agents that both hesitate (high entropy) at a
decision point. A coding agent hesitates because multiple valid
implementations exist; rollouts explore these alternatives
productively, so high entropy positively predicts rollout utility
(Type~D). A fact-verification agent hesitates because the
retrieved evidence is ambiguous; rollouts from the same model
draw observations from the \emph{same} information-bounded
environment, so they cannot synthesize evidence the environment
did not provide, and high entropy instead predicts rollout
\emph{futility} (Type~I). The same numerical signal carries
opposite meaning because the \emph{source} of uncertainty
differs.
=======
now ask \emph{why}. We model direction reversal as arising from two coexisting state types with opposite signal--utility relationships. As discussed in \S\ref{sec:intro}, high entropy carries opposite
meaning depending on whether the agent lacks information
(Type~I, where rollouts cannot help) or faces multiple viable
options (Type~D, where rollouts productively explore
alternatives).
>>>>>>> 1bf5b0e2c05afb339eb3e76f610bbf4ff034c5a3

\paragraph{Model.}
We model direction reversal as arising from two state types at each
decision step:
\begin{itemize}[leftmargin=*,nosep]
\item \textbf{Type~I (information poverty):} Utility is
  \emph{negatively} related to entropy:
  $U_I(s) \sim -\alpha H(s) + \varepsilon_I$, $\alpha > 0$.
\item \textbf{Type~D (decision difficulty):} Utility is
  \emph{positively} related to entropy:
  $U_D(s) \sim +\beta H(s) + \varepsilon_D$, $\beta > 0$.
\end{itemize}
Let $p_I(\mathcal{E})$ denote the fraction of Type~I states in
environment $\mathcal{E}$. The marginal correlation becomes
\begin{equation}
  \rho(\mathcal{E}) \;\approx\;
  \beta - (\alpha + \beta)\,p_I(\mathcal{E}),
  \label{eq:direction}
\end{equation}
which reverses sign at the critical threshold
$p_I^* = \beta/(\alpha+\beta)$. Environments with
$p_I > p_I^*$ exhibit negative $\rho$ (Type~I dominated); those
with $p_I < p_I^*$ exhibit positive $\rho$ (Type~D dominated);
near $p_I^*$, the two effects cancel and entropy carries
<<<<<<< HEAD
negligible signal. This prediction is consistent with the observed
Qwen3 pattern in Table~\ref{tab:signal-discovery}: search-based QA
(FEVER, HotpotQA, TWExpress) shows negative $\rho$, code
generation (APPS) shows positive $\rho$, and weak-signal
environments (Plancraft) cluster near zero.

\paragraph{Backbone-level reversal is a prediction, not an
anomaly.} A subtlety is essential here: $p_I$ is not a property
of the environment alone. A state is information-poor only
relative to what the agent already encodes; the same observation
can be Type~I for one model and Type~D for another. The model
backbone therefore enters the mixture directly, and the effective
$p_I$ is indexed by the (environment, backbone) pair. Read this
way, the backbone-level sign flips of
Table~\ref{tab:signal-discovery}---most extreme on FEVER between
Phi-3.5 and Llama-3.1---are exactly what the model predicts:
backbone changes shift $p_I$, and when they cross $p_I^*$, $\rho$
flips. Far from contradicting the two-source view, backbone
reversal becomes one of its predictions
(Appendix~\ref{app:multi-backbone}).
=======
negligible signal. This expression captures the sign-level
dependence of the marginal trend on the mixture proportion under
a simplified linear model; it is not intended as a precise
estimator of Spearman $\rho$. The prediction is consistent with
the observed Qwen3 pattern in Table~\ref{tab:signal-discovery}:
search-based QA (FEVER, HotpotQA, TWExpress) shows negative
$\rho$, code generation (APPS) shows positive $\rho$, and
weak-signal environments (Plancraft) cluster near zero.

\paragraph{Backbone-level reversal is a prediction, not an
anomaly.} $p_I$ is not a property of the environment alone. A
state is information-poor only relative to what the agent already
encodes; the same observation can be Type~I for one model and
Type~D for another. The model backbone therefore enters the
mixture directly, and the effective $p_I$ is indexed by the
(environment, backbone) pair. The backbone-level sign flips of
Table~\ref{tab:signal-discovery}, most extreme on FEVER between
Phi-3.5 and Llama-3.1, are exactly what the model predicts:
backbone changes shift $p_I$, and when they cross $p_I^*$, $\rho$
flips (Appendix~\ref{app:multi-backbone}).
>>>>>>> 1bf5b0e2c05afb339eb3e76f610bbf4ff034c5a3

\paragraph{Theoretical grounding.}
This direction reversal is a special case of Simpson's
paradox~\cite{simpson1951interpretation,DBLP:books/acm/22/Pearl22n}:
aggregating subpopulations with opposing within-group trends
reverses the aggregate correlation. Our decomposition adapts the
epistemic/aleatoric distinction~\cite{der2009aleatory} to the
meta-decision setting: Type~I states exhibit epistemic-like
uncertainty (information deficit), while Type~D states exhibit
aleatoric-like diversity (multiple viable paths). We treat this
two-source decomposition as a \emph{minimal falsifiable} model
rather than a post-hoc explanation: in
Section~\ref{sec:theory-verification} we derive three predictions
from it (temporal dynamics, cross-environment consistency, and
signal identity alignment) and verify each on held-out data.

<<<<<<< HEAD
The decomposition also makes the insufficiency claim of
\S\ref{sec:empirical-landscape} precise. Since the latent type
$\tau(s) \in \{I, D\}$ governs the sign of $\mathrm{VOC}(s)$
but is not measurable from $\sigma(s)$ alone, no function of
$\sigma$ can recover $\tau$, and therefore no function of
$\sigma$ can recover the sign of $\mathrm{VOC}$. The
information-theoretic content of the reversal is exactly this
non-identifiability of $\tau$ from $\sigma$.

% §3.3 (Necessity of Direction Discovery) folded into §3.1
% "Implication." paragraph above; full proposition and proof are
% deferred to Appendix~\ref{app:proof-necessity}.
=======
\paragraph{Implication.}
Under this decomposition, $\sigma$ is not a sufficient statistic
for $\mathrm{VOC}$: two states with identical $\sigma$ but
opposite latent types have opposite-sign utility. No policy of the
form $g(\sigma(s))$ can therefore deliver non-negative value across
heterogeneous settings, and uncertainty-only fixed-direction gates
empirically fail across our six environments
(Proposition~\ref{prop:necessity}, proof in
Appendix~\ref{app:proof-necessity}). Recovering sufficiency
requires augmenting $\sigma$ with features that distinguish the two
state types, directly motivating the multi-feature design of
\S\ref{sec:method}.

% \paragraph{A diagnostic for practitioners.}
% Collecting a few hundred signal-agnostic $(s, U)$ pairs and
% reporting $\rho(\sigma, U)$ with its bootstrap confidence interval
% is sufficient to decide whether a fixed direction is safe for a
% given deployment. A confidence interval that crosses zero
% indicates that fixed-direction gating is unreliable; a stable
% single-sign correlation indicates it is safe.
>>>>>>> 1bf5b0e2c05afb339eb3e76f610bbf4ff034c5a3
